\chapter{Kết luận}

\section{Trả lời các câu hỏi nghiên cứu}

Khóa luận nghiên cứu cách giới hạn việc điều phối làm thay đổi giờ đã báo cho hành khách,
đồng thời đo chi phí phục vụ và ranh giới của bảo đảm đó. Kết quả cần được đọc theo hai mức:
một bất biến của cơ chế được chứng minh từ định nghĩa và các giả thiết; hiệu quả vận hành
được đo trong những thiết lập mô phỏng cụ thể. Hai mức này không thay thế nhau.

\paragraph{RQ1: Có thể cưỡng chế chính xác lời hứa nào?}
Phép phân rã ba mốc so phương án mới với lời hứa đã công bố và với phép chiếu tuyến giữ
nguyên trên trạng thái hiện tại. Sổ cái cộng dồn phần dịch chuyển do quyết định, riêng cho
từng khách và từng chiều cam kết. Cổng loại phương án vượt trần trước khi bộ giải chọn
phương án. Với tuyến giữ nguyên, phần dịch chuyển này bằng không. Khi sổ cái trước còn
trong trần và phương án giữ tuyến qua các kiểm tra khác, cổng hành vi không loại phương án
đó; đây là bảo đảm trong Định lý~\ref{dl:hocong}.

Bảo đảm này nói về phần dịch chuyển tính so với tuyến giữ nguyên, không nói rằng giờ trả
thực tế luôn nằm trong một cửa sổ quanh lời hứa đầu. Mốc giữ tuyến còn phản ánh tiến độ xe,
thời gian thực thi và cập nhật thời gian đi. Vì vậy phép phân rã là một phép hạch toán vận
hành tại quyết định, không nhận diện tác động nhân quả thuần của giao thông. Khóa luận cũng
không coi việc kiểm soát lời hứa hay đặt hạn chót là mới: các cơ chế đó đã có trong văn liệu.
Phần nghiên cứu ở đây là đối tượng được cộng vào sổ cái và khung so sánh bốn họ cổng dưới
trôi dạt.

\paragraph{RQ2: Giá phục vụ của việc giảm dịch chuyển là bao nhiêu?}
Trên panel tĩnh chính, gói cam kết C1 giảm mạnh phần dịch chuyển do điều phối, nhưng số
khách phục vụ giảm từ $1.735$ xuống $1.581$ trên $2.160$ yêu cầu, tức $7{,}13$ điểm phần
trăm. Đây là kết quả của cả gói cấu hình, gồm khóa xác nhận và trần giờ trả; không được quy
toàn bộ chi phí cho riêng trần 30 giây. Chi phí này vượt mức một điểm phần trăm đã đặt,
nên tiêu chí tổng hợp H6 không đạt dù các kiểm tra cơ chế đạt. Thang ngân sách giữ các khóa cố định cho phép xem
phần đánh đổi khi thay trần. Những khác biệt nhỏ về phục vụ giữa hai trần có thể đổi chiều
theo ô, do mỗi lựa chọn sớm làm thay đổi trạng thái các quyết định sau.

Ở thí nghiệm T38 với thời gian đi thay đổi, $34$ ô kiểm tra cho một kết quả hẹp hơn:
cổng $C$-30 phục vụ $2.338$ khách, bằng $M^{+}$-30 và gần $V$-30 ($2.336$ khách), khi dùng
cùng luật phục hồi không làm tệ hơn tham chiếu. So với nhánh không cổng $U$, $C$-30 mất
$96$ khách trên $3.672$ yêu cầu, tức $2{,}61$ điểm phần trăm tỉ lệ phục vụ
(mục~\ref{muc:t38}). Do đó dữ liệu không ủng hộ kết luận rằng cổng hành vi phục vụ nhiều
hơn hai họ còn lại. Kết quả cho biết chi phí của mỗi cấu hình trong thiết lập đã thử,
không xác định một trần tối ưu chung cho vận hành thực tế.

\paragraph{RQ3: Trôi dạt và phục hồi thay đổi khả thi và lời hứa thế nào?}
Khi thời gian đi thay đổi, tuyến giữ nguyên vẫn có thể vượt hạn chót hoặc tiêu hết ngân
sách biến động khách nhìn thấy. Các thí nghiệm bốn họ cổng phân biệt được việc bị cổng cam
kết loại với việc dừng vì kiểm tra vật lý. Chẳng hạn, cổng hành vi không gây vô nghiệm do
cam kết trên panel trôi dạt sáng, nhưng hệ thống vẫn dừng ở $11/20$ ô vì cửa sổ đón. Do
đó không được diễn giải bất biến của cổng thành bảo đảm hệ thống luôn tiếp tục phục vụ.

Phục hồi cho phép xe tiếp tục hoạt động khi tuyến giữ nguyên đã vi phạm một cổng. Việc
tiếp tục này phải kèm bản ghi vi phạm; nó không có nghĩa lời hứa cũ đã được giữ. Trên các
ô kiểm tra T38, $C$-30 không có quyết định cần phục hồi và không vượt trần hành vi.
Trong khi số khách gần như bằng nhau, $M^{+}$-30 và $V$-30 vẫn phá lời hứa theo thước đo
của chính họ. Tỉ lệ khách bị phá lời hứa, lấy trung bình không trọng số theo ô,
lần lượt khoảng $14{,}5\%$ và $21{,}6\%$. Nếu chấm $C$ bằng thước đo hạn chót hoặc biến
động nhìn thấy, khách của $C$ cũng có thể vượt các mức ấy. Sự khác nhau nằm ở nội dung
của bảo đảm và nhu cầu phục hồi, không phải ở một trải nghiệm giờ giấc chắc chắn tốt hơn.

Khoảng dự phòng của T38 bổ sung một kết quả thực nghiệm cho biến động khách nhìn thấy.
Từ $15$ ô hiệu chỉnh, $\gamma=222{,}259$ giây tạo cận
$\beta+\gamma=252{,}259$ giây với $\beta=30$ giây. Có $32/34$ ô kiểm tra được phủ
theo điểm số trôi dạt; chỉ một ô có biến động nhìn thấy thực sự vượt cận này.
Hai đại lượng không đồng nhất: vượt điểm số dự phòng chưa chắc làm tổng biến động nhìn
thấy vượt cận. Độ phủ ở riêng khung 7:00--9:00 chỉ là $7/9$, nên kết quả chung không
cho một bảo đảm xác suất 90\% cho từng khung giờ hay cho các ngày ngoài mẫu.

\section{Ý nghĩa và giới hạn}

Bài học thiết kế là phải công bố rõ đại lượng mà hệ thống cam kết. Giới hạn phần dịch
chuyển do quyết định có thể kiểm ngay trước khi chọn phương án và giữ được bất biến
giữ tuyến. Giới hạn giờ khách nhìn thấy cần thêm mô hình biến động, khoảng dự phòng và
quy tắc xử lý khi bảo đảm bị phá. Các chỉ số phục vụ, nội dung lời hứa, vi phạm và phục
hồi phải được báo riêng; cùng số khách không đủ chứng minh hai cấu hình tương đương
về bảo đảm.

Những giới hạn ở mục~\ref{muc:hanche} ảnh hưởng trực tiếp đến cách áp dụng kết quả.
Mỗi ô chỉ có $108$ yêu cầu và $8$ xe. Mạng thử là bao đóng thời gian đi giữa $96$
điểm, với tọa độ tổng hợp, không phải mạng đường phố Manhattan đầy đủ. Nhu cầu và hệ số
giao thông có nguồn dữ liệu thực, nhưng hệ số giao thông được áp vô hướng cho toàn mạng
theo mốc 15 phút. Một hạt giống cố định và một miền dữ liệu chưa đủ để đánh giá tính
ổn định trên nhiều môi trường vận hành.

Các thử nghiệm trôi dạt được báo như bằng chứng thăm dò có niêm phong, với mức kiểm
chứng được nêu cho từng panel. T38 có tập kiểm tra tách theo nhóm ngày--mẫu, nhưng
ngưỡng và cách chia đã được chọn sau khi xem kết quả phát triển. Có 75 chuyến nguồn
chung giữa hiệu chỉnh và kiểm tra. Các khung cùng ngày vẫn liên quan; khoảng lấy lại
mẫu theo ô không loại bỏ sự phụ thuộc này. Khoảng dự phòng còn do một ô hiệu chỉnh
quyết định. Những giới
hạn ấy không xóa kết quả quan sát, nhưng ngăn suy rộng độ phủ hoặc các khoảng bất
định thành bảo đảm cho toàn bộ nhu cầu.

Bộ sinh chính giữ thứ tự khách cũ và chèn thêm điểm dừng. Trong điều kiện đó, phần
dịch chuyển do quyết định gần như chỉ theo chiều muộn, nên nhiều trường hợp chặn
chuỗi và chặn độ lệch hai đầu gần trùng nhau. Một bộ sinh biết sửa tuyến có thể tạo
kết quả khác. Hơn nữa, sổ cái trị tuyệt đối còn tính phí cho quyết định kéo giờ về
gần lời hứa đầu, kể cả khi đó là sửa chữa có lợi cho khách. Khóa luận chưa đo lựa
chọn, mức chấp nhận hay sự hài lòng của hành khách; giảm độ dịch chuyển không tự
chứng minh lợi ích kinh tế hoặc tâm lý.

\section{Hướng tiếp theo}

\textbf{Hoàn thiện kiểm chứng cho thời gian đi thay đổi.} Ưu tiên là một bộ kiểm
chứng độc lập chấp nhận nhiều mốc thời gian đi và kiểm lại cùng các hợp đồng về
lời hứa, vật lý, vi phạm và phục hồi. Việc này cần phân biệt lỗi hạ tầng, lỗi
adapter và kết cục đúng thiết kế, để một lần chạy hoàn tất không bị hiểu nhầm là
mọi bảo đảm đều đã thỏa.

\textbf{Kiểm tra khả năng chuyển kết quả sang môi trường khác.} Cần thử trên mạng
đường phố lớn hơn, dữ liệu nhiều ngày và các nhóm kiểm tra độc lập. Thiết kế phải
tính sự phụ thuộc giữa các khung cùng ngày và báo độ phủ theo khung, thay vì chỉ
gộp chung. Câu hỏi cần trả lời là khoảng dự phòng và chi phí phục vụ có còn hợp lý
khi thay cấu trúc mạng, mật độ nhu cầu và dạng biến động thời gian đi hay không.

\textbf{Đánh giá sửa chữa có lợi và ý nghĩa đối với người dùng.} Trước hết phải
thống nhất cách khóa chặng xe đang chạy giữa mô hình miền và bộ mô phỏng, rồi mới
so bộ sinh sửa tuyến với bộ sinh chèn trên cùng hợp đồng. Khi đó có thể kiểm một
sổ cái chỉ tính phần làm xấu lời hứa, thay vì mặc định mọi dịch chuyển đều có hại.
Một nghiên cứu người dùng riêng cần xác định hành khách hiểu và coi trọng loại
lời hứa nào. Hai phép kiểm này giúp quyết định nên giữ bảo đảm hành vi hiện tại,
bổ sung bảo đảm kết quả, hay thay đổi đại lượng được tính phí.

Kết luận được hỗ trợ là: sổ cái hành vi bảo vệ một đại lượng có thể cưỡng chế,
với chi phí phục vụ đo được trong mô phỏng. Nó cung cấp cách phân biệt hành vi
điều phối, thay đổi khi giữ tuyến và kết quả khách thấy; nó chưa thay thế một
bảo đảm giờ đến thực tế trên hệ vận tải quy mô lớn.
